%% Brutalist / old-PC CV template - [YOUR_NAME]
%% Monospace throughout, boxed header, inverted black section bars, hard rules.
%% This is the port's default layout. cv/main_moderncv.tex keeps the moderncv
%% banking template, whose text layer extracts in cleaner reading order - prefer
%% it when the reader is an ATS parser.
%% Compile with: lualatex main_example.tex
\documentclass[10pt]{article}

\usepackage[scale=0.80]{geometry}
\usepackage{iftex}
\ifpdftex\usepackage[T1]{fontenc}\fi
\usepackage[utf8]{inputenc}
\usepackage{lmodern}
\usepackage{xcolor}
\usepackage[hidelinks]{hyperref}
\sloppy
\setlength{\emergencystretch}{3em}

\pagestyle{empty}
\setlength{\parindent}{0pt}
\setlength{\parskip}{2pt}
\renewcommand{\familydefault}{\ttdefault}
\renewcommand{\labelitemi}{\textbullet}
\setlength{\leftmargini}{1.1em}
\setlength{\itemsep}{0.5pt}
\setlength{\topsep}{1pt}
\setlength{\parsep}{0pt}

% Inverted section bar: black band, white uppercase text, rule under it.
\newcommand{\sectionbar}[1]{%
  \par\vspace{7pt}%
  \noindent\colorbox{black}{\parbox{\dimexpr\linewidth-2\fboxsep}{\color{white}\bfseries\footnotesize\MakeUppercase{#1}}}%
  \par\vspace{2pt}%
  \noindent\rule{\linewidth}{0.8pt}%
  \par\vspace{2pt}%
}

% Entry header: role on the left, dates hard-right, thin rule under it.
\newcommand{\entry}[3]{%
  \par\vspace{2pt}%
  {\bfseries #1}\hfill{\footnotesize #2}\par
  {\footnotesize #3}\par\vspace{1pt}%
  \noindent\rule{\linewidth}{0.3pt}\par\vspace{2pt}%
}

\begin{document}

% ---------- HEADER ----------
\noindent\fbox{\parbox{\dimexpr\linewidth-2\fboxsep-2\fboxrule}{%
  {\LARGE\bfseries [YOUR NAME]}\par\vspace{3pt}%
  {\small [YOUR ROLE]}\par\vspace{4pt}%
  \rule{\linewidth}{0.4pt}\par\vspace{3pt}%
  {\footnotesize
   [Your Address, City, Country] \textbar{} [+XX XXXXXXXXXX] \textbar{} [your.email@example.com]\par\vspace{1pt}%
   github.com/[your-username] \textbar{} linkedin.com/in/[your-profile]\par\vspace{1pt}%
   [Work authorization] \textbar{} [remote, hybrid or relocation]}
}}

\vspace{4pt}
{\footnotesize
[PROFILE STATEMENT - 3 to 5 lines, tailored per role. What you do, on what stack, what you have
shipped. Facts and effects only: no adjectives, no business rules of a client.]}

% ---------- COMPETENCIES ----------
\sectionbar{Core competencies}

{\footnotesize
\textbf{Languages:} [languages].\par
\textbf{Backend:} [frameworks, runtimes, data access, API design, auth].\par
\textbf{Frontend:} [frameworks, versions, state, styling, mobile].\par
\textbf{Data and testing:} [databases, data tools; test frameworks].\par
\textbf{Other:} [game dev, audio, tooling, AI tooling, deployment].}

% ---------- AGENTIC ----------
\sectionbar{Agentic development}

{\footnotesize [How you use agent runtimes: which ones, what you author (skills, entry points),
how you drive them through changes and how you review their output before it ships.]}

% ---------- EXPERIENCE ----------
\sectionbar{Professional experience}

\entry{[Job Title]}{[YYYY -- YYYY]}{[Company] \textbar{} [location or remote] \textbar{} [contract type, team size]}

{\footnotesize
\begin{itemize}
  \item {Scope: what you work on, decisions you take, review, contact. Volume counts if measured.}
  \item {Capability 1: what you build, with the stack named.}
  \item {Capability 2: integrations, fallbacks, error handling.}
  \item {Capability 3: performance work. State the effect, never the mechanism.}
  \item {Quality: tests, review, documentation.}
\end{itemize}}

% ---------- EDUCATION ----------
\sectionbar{Education}

\entry{[Degree or programme]}{[YYYY -- YYYY]}{[Institution] \textbar{} [city, country]}
{\footnotesize [Official title on completion, certification status, track.]}

\vspace{2pt}
\entry{[Secondary title]}{[YYYY -- YYYY]}{[School] \textbar{} [city]}
{\footnotesize [Orientation or specialisation.]}

% ---------- PROJECTS ----------
\sectionbar{Projects}

{\footnotesize
\textbf{[Project name]} ([own project / team project / public], [status]) --- [what it is, how it is
built, what it proves].\par\vspace{2pt}
\textbf{[Project name]} ([own project], [status]) --- [what it is, architecture, tests].}

% ---------- LANGUAGES ----------
\sectionbar{Languages}

{\footnotesize [Language] ([level]); [Language] ([level]).}

% ---------- REFERENCES ----------
\sectionbar{References}

{\footnotesize Available on request.}

\end{document}
